\par\noindent\textbf{(1)} 珠子位置在水平面到轴距离 \(R\sin\theta\)，速度平方 \(v^2=R^2\dot\theta^2+\Omega^2R^2\sin^2\theta\)。因此 \(L=\tfrac12mR^2\dot\theta^2+\tfrac12m\Omega^2R^2\sin^2\theta-mgR(1-\cos\theta)\)，其有效势是 \(U=m gR(1-\cos\theta)-m\Omega^2R^2\sin^2\theta/2\)。
\par\noindent\textbf{(2)} 平衡条件 \(U'=mR\sin\theta(g-\Omega^2R\cos\theta)=0\)。底点 \(\theta=0\) 总是平衡；当 \(\Omega^2R>g\) 时还出现两个对称平衡 \(\theta=\pm\arccos(g/(\Omega^2R))\)，临界 \(\Omega_c=\sqrt{g/R}\)。
\par\noindent\textbf{(3)} 底点 \(U''(0)=mR(g-\Omega^2R)\)，所以 \(\omega_0^2=g/R-\Omega^2\)。对非零平衡使用 \(\cos\theta_*=g/(\Omega^2R)\)，可得 \(U''(\theta_*)=m\Omega^2R^2\sin^2\theta_*\)，故 \(\omega_*^2=\Omega^2[1-g^2/(\Omega^4R^2)]>0\)。
\par\noindent\textbf{(4)} 在 \(\theta\approx0\) 附近展开 \(U=\tfrac12mR(g-\Omega^2R)\theta^2+\frac{mR}{24}(-g+4\Omega^2R)\theta^4+O(\theta^6)\)。在阈值 \(\Omega^2R=g\) 时四次系数 \(m gR/8>0\)，故平衡连续分裂为超临界叉形分岔。
